% !TeX root = ../proyecto.tex

\begin{UseCase}{CU-03}{Investigación profunda bajo demanda}{
	Permite al analista activar el módulo DeepInvestigator sobre una noticia clasificada como ``Alta'' para generar una ficha de caso estructurada mediante búsqueda web y síntesis con IA (Gemini Flash).
}
	\UCitem{Versión}{\color{Gray}
		1.0
	}\UCitem{Autor}{\color{Gray}
		Héctor Alberto Morales Martínez
	}\UCitem{Supervisa}{\color{Gray}
		Director de Trabajo Terminal
	}\UCitem{Actor}{
		Analista de datos / Trabajador social
	}\UCitem{Propósito}{\begin{Titemize}
		\Titem Profundizar en los detalles de un caso específico delegando la investigación multi-fuente y la síntesis a agentes de IA.
	\end{Titemize}
	}\UCitem{Entradas}{\begin{Titemize}
		\Titem ID de la noticia seleccionada.
	\end{Titemize}
	}\UCitem{Origen}{\begin{Titemize}
		\Titem Botón ``Investigar mas en la web'' en la tarjeta de la noticia dentro del Dashboard.
	\end{Titemize}
	}\UCitem{Salidas}{\begin{Titemize}
		\Titem Ficha estructurada en formato JSON renderizada en la interfaz (Ubicación, Víctimas, NNA Afectados, Situación Actual, Medios de Contacto).
	\end{Titemize}
	}\UCitem{Destino}{\begin{Titemize}
		\Titem Interfaz del Dashboard (Casos objetivo).
		\Titem Base de datos (almacenamiento de la investigación asociada a la noticia).
	\end{Titemize}
	}\UCitem{Precondiciones}{\begin{Titemize}
		\Titem La noticia debe estar registrada en el sistema y tener clasificación ``Alta''.
		\Titem Debe existir conexión a Internet y cuota disponible en la API de Gemini.
	\end{Titemize}
	}\UCitem{Postcondiciones}{\begin{Titemize}
		\Titem Se vincula una ficha de investigación estructurada al registro de la noticia.
	\end{Titemize}
	}\UCitem{Errores}{\begin{Titemize}
		\Titem {\bf \hypertarget{CU-03.E1}{E1}}: DuckDuckGo no retorna resultados válidos o bloquea la conexión \Trayref{CU-03}{A}.
		\Titem {\bf \hypertarget{CU-03.E2}{E2}}: Exceso de cuota en la API de Gemini o fallo de red con el proveedor \Trayref{CU-03}{B}.
	\end{Titemize}
	}\UCitem{Tipo}{
		Caso de uso primario
	}\UCitem{Observaciones}{
		La operación se ejecuta de manera asíncrona en un hilo de fondo (background thread) para no bloquear la interfaz web del usuario.
	}
\end{UseCase}

%--------------------------------------
\begin{UCtrayectoria}
	\UCpaso[\UCActor] Analiza el listado de noticias con score ``Alta'' en el dashboard.
	\UCpaso[\UCActor] Selecciona una noticia de interés y hace clic en el botón \IUbutton{Investigar mas en la web}.
	\UCpaso[\UCsist] Cambia el estado de la investigación a \texttt{running} y lanza un hilo de ejecución en segundo plano.
	\UCpaso[\UCsist] Invoca a la API de Gemini para generar una consulta de búsqueda altamente específica (máximo 6 palabras) basada en la noticia \ErrorRef{CU-03}{E2}{Fallo API Gemini}.
	\UCpaso[\UCsist] Ejecuta la búsqueda de la consulta en DuckDuckGo (HTML y Lite) \ErrorRef{CU-03}{E1}{Sin resultados DDG}.
	\UCpaso[\UCsist] Extrae y limpia las URLs resultantes de la búsqueda.
	\UCpaso[\UCsist] Realiza scraping del contenido de las fuentes complementarias utilizando \texttt{StealthSession}.
	\UCpaso[\UCsist] Envía el texto multi-fuente consolidado a la API de Gemini junto con el prompt de analista criminal \ErrorRef{CU-03}{E2}{Fallo API Gemini}.
	\UCpaso[\UCsist] Recibe y valida el JSON estructurado con los detalles del caso.
	\UCpaso[\UCsist] Guarda el resultado en la base de datos y actualiza el estado a \texttt{done}.
	\UCpaso[\UCActor] El frontend, mediante \textit{polling}, detecta el fin del proceso y despliega la ficha estructurada en pantalla.
\end{UCtrayectoria}

%--------------------------------------
\begin{UCtrayectoriaA}{CU-03}{A}{DuckDuckGo no retorna URLs utilizables o bloquea la petición}
	\UCpaso[\UCsist] Activa el mecanismo de ''fallback'' utilizando la herramienta nativa \texttt{GoogleSearch} de Gemini.
	\UCpaso[\UCsist] Ejecuta la búsqueda de información directamente a través del modelo de lenguaje sin scraping externo.
	\UCpaso[\UCsist] Si la búsqueda interna también falla, marca el estado como \texttt{error} y notifica al usuario en el dashboard.
	\UCpaso[] El Caso de Uso termina o continúa en el paso 8 de la trayectoria principal si el fallback es exitoso.
\end{UCtrayectoriaA}

%--------------------------------------
\begin{UCtrayectoriaA}{CU-03}{B}{La API de Gemini responde con un error de cuota o falla de conexión}
	\UCpaso[\UCsist] Captura la excepción de la API externa.
	\UCpaso[\UCsist] Registra el error en los logs.
	\UCpaso[\UCsist] Cambia el estado de la investigación a \texttt{error} indicando ``Fallo en el servicio''.
	\UCpaso[\UCActor] Visualiza el mensaje de error en la interfaz, permitiéndole reintentar más tarde.
	\UCpaso[] El Caso de Uso termina de manera anormal.
\end{UCtrayectoriaA}
